\cotefile{156-8}
\scanpage{50}
\mspage{41}
% archive sheet number: 50

$(\mathrm{i})\Longrightarrow(\mathrm{ii})$ Je \unsure{note} : prouvons la formule
\[
X_i\cap X_j=\bigcup_{k\leq i,j}X_k
\]
(côté $\supset$ est évident). Mais par $3'$ on a
\[
X_i=\bigcup_{i'\leq i}X_{i'}^\circ,\quad X_j=\bigcup_{j'\leq j}X_{j'}^\circ
\]
d'où
\[
X_i\cap X_j=\bigcup_{\substack{i'\leq i\\j'\leq j}}X_{i'}^\circ\cap X_{j'}^\circ
\]
Or par $2')$ on a $X_{i'}^\circ\cap X_{j'}^\circ=\varnothing$ sauf si $i'\ne j'$, d'où
\[
X_i\cap X_j=\bigcup_{k\leq i,j}X_k^\circ\subset\bigcup_{k\leq i,j}X_k\qquad\text{qfd}
\]

$(\mathrm{ii})\Longrightarrow(\mathrm{i})$. Cela résulte du

\deleted{\uline{Corollaire}} \inserted{Lemme} Supposons qu'on ait, $\forall i,j\in I$
\[
X_i\cap X_j\supset\bigcup_{k\leq i,j}X_k.
\]
Alors on a les propriétés :
\[
\forall i,j\in I,\ i\ne j\quad\deleted{\unsure{X_j}}\quad X_i^\circ\cap X_j^\circ=0\tag*{$2')$}
\]
\inserted{Si de plus $\bigcup_{i\in I}X_i^\circ=|\unsure{\Phi|P}|$ on a aussi}
\[
\forall i\in I\qquad X_i=\bigcup_{j\leq i}X_j^\circ.\tag*{\unsure{$3')$}}
\]

Prouvons $2')$ on a \inserted{\unsure{car|or}}
\[
X_i^\circ\cap X_j^\circ\subset X_i\cap X_j=\unsure{\bigcup|\coprod}_{k\leq i,j}X_k
\]
et le deuxième \unsure{mbr}, si $i\ne j$, \deleted{\unsure{et}} tout $k\leq i,j$ est ou bien $<i$, ou bien $<j$ (il ne peut être \unsure{égal} à la fois : $i$ et $j$) donc \deleted{\unsure{on}} $X_k\subset\partial X_i$ ou $X_k\subset\partial X_j$, donc
\[
i\ne j\Longrightarrow X_i\cap X_j\subset\partial X_i\cup\partial X_j\subset\complement(X_i^\circ\cap X_j^\circ)
\]
donc
\[
X_i^\circ\cap X_j^\circ\subset\complement X_i^\circ\cap X_j^\circ\qquad\text{d'où}\quad X_i^\circ\cap X_j^\circ=\varnothing
\]

Prouvons $3')$. \inserted{On a bien sûr $\supset$, il faut voir $\subset$.} Soit $x\in X_i$, par hyp. $\exists j\in I$ tel que $x\in X_j^\circ$, \unsure{tout} revient à voir que $X_j^\circ\subset X_i$.
